\section{Application: City-Scale Source Attribution}
\label{app:application}

This appendix gives the construction of the operator, the nuisance basis and the
platform for the pollution application summarized in Section~\ref{sec:problem_setup}.
In the notation of the general model, \(A(\omega)=H_\Phi^{\mathrm{lag}}\),
\(\widetilde A=\widetilde H_\Phi\), the columns \(\mathbf h_{kb}\) are the
fingerprints \(\mathbf a_j\) with \(j=(k,b)\), and the background basis \(Q\) is the
nuisance basis.

\subsection{Transport Model, Inventories and Lagged Response}
\label{app:transport_model}


\noindent\textbf{City-scale transport model.}
Let \(u(\mathbf x,t)\) be the pollutant concentration field on
\(\Omega\subset \mathbb R^d\), evolving as \(\partial u/\partial t
= \mathcal F_\eta(u,\mathbf x,t;w(t)) + s(\mathbf x,t)+b(\mathbf x,t)\) under
meteorological state \(w(t)\), local emission \(s\), background \(b\), and
transport/dispersion parameters \(\eta\); the dynamics may be realized by a plume
model, an advection--diffusion solver, or a transport-constrained surrogate.

\noindent\textbf{Sparse observations.}
Discretize the concentration field on \(n\) grid cells, so that
\(\mathbf u_t\in \mathbb R^n\).  Let
\(X_{\mathrm{sens}}=\{\mathbf x_1,\ldots,\mathbf x_m\}\) be the fixed sensor
locations with \(m\ll n\).  The observation operator
\(O:\mathbb R^n\to \mathbb R^m\) samples or interpolates the latent field at
these locations, \(\mathbf y_t = O\mathbf u_t + \boldsymbol\epsilon_t\); stacking
over \(t=1,\ldots,T\) gives
\(Y = [\mathbf y_1^\top,\ldots,\mathbf y_T^\top]^\top \in \mathbb R^{mT}\).
Because \(m\) is small relative to the field dimension, many latent fields and
source configurations fit the same sensor~trajectory.

\noindent\textbf{Inventory-based source parameterization.}
We assume a set of \(K\) known or proxy source maps,
\(S = [\mathbf s_1\;\cdots\;\mathbf s_K]\in \mathbb R^{q\times K}\),
where \(\mathbf s_k\in\mathbb R^q\) is the spatial pattern of source group
\(k\).  Time-varying source activity uses a fixed, predeclared nonnegative
temporal dictionary
\(\Phi=[\boldsymbol\phi_1\;\cdots\;\boldsymbol\phi_B]\in\mathbb R_+^{T\times B}\)
and coefficients \(C=[c_{kb}]\in\mathbb R_+^{K\times B}\), so the activity of
source \(k\) is \(\theta_k(t)=\sum_{b=1}^{B}c_{kb}\phi_b(t)\ge0\).  The basis is a
shared dictionary of admissible temporal patterns, not a claim that all sources
follow one trajectory: each source has its own nonnegative weights, and
source-specific knowledge is imposed by masking inadmissible coefficients.  We do
not learn \(\Phi\) from the same sparse observations, since jointly estimating
\(\Phi\) and \(C\) would add bilinear non-identifiability beyond the source
distinguishability analyzed here.  We vectorize \(C\) in source-major, basis-minor order as
\(\mathbf c=\operatorname{vec}_{\mathrm{src}}(C)\in\mathbb R_+^J\), \(J=KB\); the
constant-activity model is the special case \(B=1\), \(\phi_1(t)=1\).

\noindent\textbf{Wind-conditioned lagged response.}\label{subsec:wind_field_estimation}
Government wind direction records where wind comes from, whereas transport needs
the direction pollution moves; we convert direction \(\mathit{WD}\) and speed
\(\mathit{WS}\) to transport vectors
\begin{equation}
(U_x,V_y)=-\mathit{WS}\,(\sin\alpha,\cos\alpha),\;\;
\alpha=\mathit{WD}\,\pi/180,
\label{eq:wind_direction_conversion}
\end{equation}
and complete them on the response grid with a normalized Gaussian-kernel
coordinate-query imputer to obtain a gridded transport field
\(\widehat{\mathbf w}_t(\mathbf x_g)\); and the same interface
supplies synthetic controlled winds and transport ensembles kept separate from
inventory scenarios.  Emissions released at \(t-\ell\) affect sensors at \(t\)
only after transport by the intervening wind.  Let
\(G_{t,t-\ell}^{\partial}(w_{t-\ell:t})\in \mathbb R^{n\times q}\) map emissions
released on the source grid at \(t-\ell\) to the concentration grid at \(t\); the
superscript \(\partial\) denotes an open-boundary operator (mass leaving
\(\Omega\) is removed, not reflected or wrapped).\label{subsec:plume_response} It
is realized by a Gaussian puff approximation evaluated at sensor coordinates that
preserves nonnegativity and is mass non-increasing inside the domain, with
dispersion, exit bookkeeping, and the pre-fit initial-condition baseline given in
Appendix~\ref{app:transport_response}.  For a maximum lag \(L\),
\(\mathcal L_t=\{0,\ldots,\min(L,t-1)\}\), the unit-coefficient sensor fingerprint
of source \(k\) and basis component \(b\) at time \(t\) is
\begin{equation}
\mathbf h_{t,kb}^{\mathrm{lag}}
= \sum_{\ell\in\mathcal L_t}
\phi_b(t-\ell)\,O G_{t,t-\ell}^{\partial}(w_{t-\ell:t})\mathbf s_k
\in \mathbb R^m.
\label{eq:constructed_response}
\end{equation}
Stacking over time in source-major, basis-minor order gives
\(H_{\Phi}^{\mathrm{lag}}\in \mathbb R^{mT\times J}\), \(J=KB\), each column the
finite-horizon sensor-time pattern produced by one unit of a coefficient.  The lag
window adds no sensor--time rows and is declared before fitting (\(L=12\) hours in every
reported run except the lag sweep of Experiment 7).  A convergence diagnostic compares
adjacent candidates \(L\), \(L+\Delta\) with the
same rows and columns,
\begin{equation}
\eta_L=
\frac{\|H_{\Phi}^{\mathrm{lag}}(L+\Delta)
-H_{\Phi}^{\mathrm{lag}}(L)\|_F}
{\max\{\|H_{\Phi}^{\mathrm{lag}}(L+\Delta)\|_F,\epsilon\}},
\label{eq:lag_convergence}
\end{equation}
and would select the smallest candidate with \(\eta_L\le\tau_L\) (default
\(10^{-3}\)); on the grid of Experiment 7 no candidate meets it.  The fitted
\(\mathbf c\) never selects \(L\).

\noindent\textbf{Background correction and projection.}\label{subsec:background_construction}
Observed concentrations include components not explained by the local inventory
(regional background, smooth trends, sensor offsets), giving
\(Y = H_{\Phi}^{\mathrm{lag}}\mathbf c + Q\boldsymbol\gamma + E\) with a
low-dimensional basis \(Q\in \mathbb R^{mT\times r}\) specified before fitting from
source-independent metadata only (timestamps, day labels, sensor identity, and
coordinates; effective rank capped at eight; rank four in every reported run except
the nuisance-basis variations of Experiments 3 and 11),
so apportionment is not confounded by background correction; columns built from
inventories, fingerprints, or fitted signals are excluded and permitted only in
labeled stress tests (Appendix~\ref{app:transport_response}).  Real records need
not be complete: let \(\mathcal O\) be the ordered observed sensor--time
PM\(_{2.5}\) rows and \(M_{\mathcal O}\in\{0,1\}^{N\times mT}\) select them, applied
identically to every row-aligned object,
\begin{equation}
Y_{\mathcal O}=M_{\mathcal O}Y,\quad
H_{\Phi,\mathcal O}^{\mathrm{lag}}
=M_{\mathcal O}H_{\Phi}^{\mathrm{lag}},\quad
Q_{\mathcal O}=M_{\mathcal O}Q,
\label{eq:observation_selection_app}
\end{equation}
with an alignment mismatch treated as an error;
wind may be imputed, PM\(_{2.5}\) is not, and \(M_{\mathcal O}=I\) for complete
controlled data.  Dropping the \(\mathcal O\) subscript, let
\(P_Q^\perp = I - QQ^\dagger\) project onto the orthogonal complement of the
background space (applied implicitly through a thin SVD of \(Q\)); the corrected
inverse problem is
\begin{equation}
\widetilde Y = \widetilde H_\Phi\mathbf c + \widetilde E,
\qquad
\widetilde Y=P_Q^\perp Y,
\qquad
\widetilde H_\Phi=P_Q^\perp H_{\Phi}^{\mathrm{lag}},
\label{eq:projected_model_app}
\end{equation}
so each projected fingerprint
\(\widetilde{\mathbf h}_{kb}=P_Q^\perp\mathbf h_{kb}^{\mathrm{lag}}\) encodes the
observable effect of one source--basis coefficient after transport, lag, sparse
sensing, and background correction.  All identifiability diagnostics in this paper
are computed from \(\widetilde H_\Phi\), not the raw inventory or the
\mbox{unprojected response}.

\subsection{Temporal-Basis Notation and Diagnostic Conventions}
\label{app:temporal_basis_notation}

\begin{table}[h]
\centering
\small
\caption{Indices and objects in the temporal-basis apportionment model.}
\label{tab:temporal_basis_notation}
\begin{tabularx}{\columnwidth}{llX}
\toprule
Symbol & Shape or range & Meaning \\
\midrule
\(k\) & \(1,\ldots,K\) & source-group index \\
\(b\) & \(1,\ldots,B\) & temporal-basis index \\
\(j=(k,b)\) & \(1,\ldots,J\) & source-major, basis-minor coefficient index \\
\(\Phi\) & \(T\times B\) & nonnegative source-activity basis \\
\(C\) & \(K\times B\) & nonnegative source--basis coefficients \\
\(\mathbf c\) & \(J=KB\) & source-major vectorization of \(C\) \\
\(H_\Phi^{\mathrm{lag}}\) & \(N\times J\) & observed-row coefficient fingerprints \\
\(\widetilde H_\Phi\) & \(N\times J\) & background-projected fingerprints \\
\(\widetilde{\mathbf h}_{kb}\) & \(N\) & fingerprint for coefficient \(c_{kb}\) \\
\bottomrule
\end{tabularx}
\end{table}

The constant-activity model uses \(B=1\), \(\phi_1(t)=1\), and
\(c_{k1}=\theta_k\).  In the general model, diagnostics are computed for the
\(J\) coefficient fingerprints.  Source-level activity trajectories are
reconstructed as
\[
\widehat\theta_k(t)=\sum_{b=1}^{B}\widehat c_{kb}\phi_b(t);
\]
fingerprint columns are not summed to create an artificial source-level
diagnostic vector.

For a source-level ambiguity summary, IASA reports the largest eligible
coherence between coefficients belonging to different sources and retains the
source--basis pair that attains it.  Coherence and ray distance are undefined
for pairs involving a coefficient with \(v_j\le\tau_v\); tables display these
entries as N/A, exclude them from pairwise extrema, and report the weak
coefficient separately.


\subsection{Transport Response Construction and Fit Details}
\label{app:transport_response}

This appendix gives the construction, background, and estimation details deferred
from Section~\ref{sec:problem_setup}.

\noindent\textbf{Wind imputer query and unit conversion.}
The observed station vectors and masks are supplied to a coordinate-query wind
imputer: given sparse station observations
\(\{(\mathbf z_i,t,\mathbf u_{i,t},m_{i,t})\}\), it predicts transport vectors
at arbitrary spatial query locations and times.  For the IASA response operator,
we query the imputer on every response-grid cell \(\mathbf x_g\) and hour
\(t\), producing
\begin{equation}
\widehat{\mathbf w}_t(\mathbf x_g)
=
\widehat f_{\omega}\!\left(\mathbf x_g,t;
\{(\mathbf z_i,t',\mathbf u_{i,t'},m_{i,t'})\}\right),
\; g=1,\ldots,n.
\label{eq:fieldformer_wind_field}
\end{equation}
This yields a gridded wind field
\(\widehat W\in\mathbb R^{T\times n\times 2}\), rather than a station-only or
city-averaged sequence.  Real-data validation is performed by masking observed
station vectors and measuring held-out station-time error; dense city-wide wind
truth is not assumed.  Meteorological speed must be converted to grid
displacement before puff advection.  For response timestep \(\Delta t_{\mathrm{s}}\)
seconds and cell dimensions \(\Delta x_{\mathrm{m}},\Delta y_{\mathrm{m}}\), we use
\begin{equation}
v_x^{\mathrm{grid}}(\mathbf x_g,t)
=
\frac{\widehat U_x(\mathbf x_g,t)\Delta t_{\mathrm{s}}}{\Delta x_{\mathrm{m}}},
\qquad
v_y^{\mathrm{grid}}(\mathbf x_g,t)
=
\frac{\widehat V_y(\mathbf x_g,t)\Delta t_{\mathrm{s}}}{\Delta y_{\mathrm{m}}}.
\label{eq:wind_unit_conversion}
\end{equation}
Every response artifact records these scales, the coordinate convention, the
wind units, and the exact converted sequence.

\noindent\textbf{Wind-field ensembles.}
The imputed gridded wind field is an estimated transport input, not ground truth.
To propagate wind-imputation uncertainty, we construct wind-field ensembles
\(\{\widehat{\mathbf w}^{(r)}_t(\mathbf x_g)\}_{r=1}^R\) using held-out-calibrated
prediction error, station bootstrap resampling, checkpoint ensembles, or
perturbations matched to validation residuals.  Each ensemble member is queried
on the same response grid and converted to grid displacement before response
construction.  The resulting response matrices are tagged as transport
ensembles and are kept separate from inventory-robustness scenarios.

\noindent\textbf{Puff dynamics and dispersion.}
We instantiate the response with a Gaussian puff approximation.  For a
unit release from source cell \(i\) at location \(\mathbf r_i\) and release time
\(\tau\), the puff center is initialized as \(\mathbf z_i(\tau)=\mathbf r_i\) and
advected by the interpolated wind field,
\(d\mathbf z_i(a)/da=\widehat{\mathbf w}_a(\mathbf z_i(a))\), \(a\in[\tau,t]\).
Source cells have integer centers and domain extents
\([-0.5,N_x-0.5]\times[-0.5,N_y-0.5]\).  With substep size \(\delta a\), the
trajectory is discretized as
\(\mathbf z_i^{r+1}=\mathbf z_i^r+\delta a\,\widehat{\mathbf w}_{a_r}(\mathbf z_i^r)\).
The final substep is shortened so that the trajectory lands exactly at the
requested observation time.  If the center exits \(\Omega\), the whole
remaining puff is removed and never reflected, wrapped, clamped, or reinserted.
Otherwise its contribution is spread with an anisotropic Gaussian kernel
\begin{equation}
K_\psi(\mathbf x;\mathbf z_i,\Sigma_i)
=\frac{\exp\left(-\frac12\|\mathbf x-\mathbf z_i\|_{\Sigma_i^{-1}}^2\right)}
{2\pi |\Sigma_i|^{1/2}},
\label{eq:gaussian_kernel}
\end{equation}
where \(\psi\) denotes dispersion parameters.  A simple covariance model is
\begin{equation}
\Sigma_i(t,\tau)
= R_i(t,\tau)
\begin{bmatrix}
\sigma_\parallel^2(t-\tau) & 0\\
0 & \sigma_\perp^2(t-\tau)
\end{bmatrix}
R_i(t,\tau)^\top,
\label{eq:dispersion_covariance}
\end{equation}
with \(R_i\) aligned to the mean wind direction along the trajectory.  When the
mean wind norm is below the numerical threshold, the downwind direction is
undefined; the implementation uses the positive-\(x\) axis as a deterministic
tie-breaker for the anisotropic covariance orientation, recorded in the response
metadata and affecting only near-calm releases.  The age in this covariance is
lower-bounded by a positive minimum dispersion time, so same-time releases remain
finite: \(\text{effective age} = \max(t-\tau, t_{\min})\).

\noindent\textbf{Sensor evaluation and loss bookkeeping.}
Response-matrix entries are computed by evaluating the kernel directly at
sensor coordinates:
\begin{equation}
\left[O G_{t,\tau}^{\partial}
(\widehat{\mathbf w}_{\tau:t};\psi)\mathbf e_i\right]_s
=\chi_i(t,\tau)K_\psi(\mathbf x_s;\mathbf z_i(t),\Sigma_i(t,\tau)).
\label{eq:sensor_response}
\end{equation}
Here \(\chi_i(t,\tau)\) is an open-boundary survival indicator for a release
from source cell \(i\) at release time \(\tau\): it is one while the advected
kernel center remains inside the modeled domain and zero after the release has
exited.  Boundary and truncation losses are recorded only as diagnostic
quantities and are never used to renormalize the sensor response.  Boundary
loss denotes the portion of the Gaussian kernel lying outside the modeled
spatial domain, while truncation loss denotes the portion omitted by evaluating
retained mass only over a finite kernel support; both are estimated separately
by grid quadrature over the response grid.  These kernel losses are distinct
from whole-release exit loss, where the advected release center leaves the open
domain and no later sensor contribution is assigned to that release.  Because
the response matrix evaluates concentration at sensor points rather than
integrating over area-weighted grid cells, summing entries across sensors or
repeated sensor-time observations does not estimate retained pollutant mass.

\noindent\textbf{Baseline policy and differentiability.}
Response matrices are constructed for source-induced increments relative to a
declared initial-condition baseline.  The baseline policy is fixed before
source fitting, recorded with the response artifact, and applied consistently to
the observations and response rows.  Application-specific choices, such as
whether to subtract a zero-source transported initial field or to use a
zero-initial-state response, are part of the experimental instantiation rather
than the generic response definition.  The open-boundary lagged plume response is
implemented in PyTorch; tensor operations retained in the graph can be
differentiated, but discrete exit events and reporting logic are not assigned
gradients.

\noindent\textbf{Background basis and implicit projection.}
The normal background basis may contain a global constant, centered and scaled
elapsed-time polynomials, local-clock daily sine/cosine harmonics,
reference-coded day effects, standardized regional sensor-coordinate trends,
reference-coded sensor offsets, and a row-aligned user basis.  It depends only
on timestamps, day labels, sensor identity, and sensor coordinates, and its
effective rank is capped at eight.  Once the primary \(Q\) is fixed, its
coefficients are estimated implicitly by projection, or explicitly in the joint
fit; each alternative source-independent basis reports the resulting visibility,
absorption, rank, conditioning, and ambiguity components.  Letting
\(Q=U\Sigma V^\top\) be a thin singular value decomposition and retaining
columns \(U_r\) above the recorded rank tolerance, projection is applied
implicitly,
\begin{equation}
P_Q^\perp X=X-U_r(U_r^\top X),
\label{eq:implicit_background_projection}
\end{equation}
without constructing an \(N\times N\) matrix.  The row metadata for \(Y\),
\(H_\Phi^{\mathrm{lag}}\), and \(Q\) must be identical: each row must refer to the
same timestamp, sensor identity, and sensor index in the same order before
projection is applied.

\noindent\textbf{Regularizers and joint fit.}
The implementation also accepts a penalty \(\lambda R(\mathbf c)\) added to the
objective of Equation~\ref{eq:nnls_estimator}, with ridge stabilization
\(R(\mathbf c)=\|\mathbf c\|_2^2\) or a weak prior around inventory-derived activity
levels \(R(\mathbf c)=\|\mathbf c-\mathbf c_0\|_2^2\); every reported run uses
\(\lambda=0\).  Equivalently, one may fit background and source coefficients jointly,
\begin{equation}
(\widehat{\mathbf c},\widehat{\boldsymbol\gamma})
=\arg\min_{\mathbf c\ge 0,\boldsymbol\gamma}
\|Y-H_{\Phi}^{\mathrm{lag}}\mathbf c-Q\boldsymbol\gamma\|_2^2
+\lambda R(\mathbf c).
\label{eq:joint_fit}
\end{equation}
The projected formulation is used for identifiability because it exposes the
part of each source--basis fingerprint that remains after background correction.

\noindent\textbf{Constrained end-to-end refinement.}
No reported run uses the refinement described here.  IASA first constructs a fixed
response matrix from the wind field and
default dispersion parameters, fits source activities, and computes
identifiability diagnostics for that declared response.  It then optionally
performs a constrained local refinement of the wind field, dispersion
parameters, source coefficients, and background coefficients, initialized at the
fixed-response solution and allowed only to make small physically constrained
corrections, so that improved sensor fit does not come from arbitrary changes to
transport.  Let \(\phi\) denote the wind-imputer parameters and \(\psi\) the
dispersion parameters of the puff response (such as \(\sigma_\parallel\),
\(\sigma_\perp\), and the minimum dispersion age), with \(\phi_0\) and \(\psi_0\)
their pretrained or default values.  We use \(R_{\mathrm{sm}}\) for a predeclared
wind-field smoothness penalty (temporal or spatial squared differences) and
\(\Psi_{\mathrm{phys}}\) for the physically admissible dispersion-parameter set.
The refinement objective is
\begin{align}
\mathcal L_{\mathrm{refine}}
=&\|Y-H_{\Phi}^{\mathrm{lag}}(\phi,\psi)\mathbf c-Q\boldsymbol\gamma\|_2^2
+\lambda_\theta R(\mathbf c) \nonumber\\
&+\lambda_w\|\widehat{\mathbf w}_{1:T}^{\phi}-\widehat{\mathbf w}_{1:T}^{\phi_0}\|_2^2
+\lambda_\psi\|\psi-\psi_0\|_2^2 \nonumber\\
&+\lambda_{\mathrm{sm}}R_{\mathrm{sm}}(\widehat{\mathbf w}^{\phi}),
\label{eq:refinement_objective}
\end{align}
subject to \(\mathbf c\ge 0\), \(\psi\in\Psi_{\mathrm{phys}}\), and
\(\|\widehat{\mathbf w}_{1:T}^{\phi}-\widehat{\mathbf w}_{1:T}^{\phi_0}\|_\infty
\le \epsilon_w\).  The refined response is accepted only as a constrained local
correction (acceptance criteria in Appendix~\ref{app:reporting_details}) and is
reported with its own response diagnostics.  Without these restrictions, wind and
plume parameters could fit sparse pollution sensors while weakening the physical
and source-level interpretation of \(\widehat{\mathbf c}\) and the reconstructed
activities.


\subsection{Response-Matrix Perturbations}
\label{app:operator_error_details}

This appendix expands the response-error discussion of
Section~\ref{sec:identifiability_theory}.  Write the projected
response discrepancy as
\(\Delta_H=\Delta_{\mathrm{tr}}+\Delta_{\mathrm{inv}}+\Delta_{\mathrm{int}}\),
where \(\Delta_{\mathrm{tr}}\) is the transport term, \(\Delta_{\mathrm{inv}}\)
the inventory term, and \(\Delta_{\mathrm{int}}\) their interaction.  The
decomposition keeps the two error sources apart: \(\Delta_{\mathrm{tr}}\) is
propagated through transport ensembles and may become an uncertainty interval,
whereas \(\Delta_{\mathrm{inv}}\) is reported only as named robustness scenarios
and is never pooled into an interval, and \(\Delta_{\mathrm{int}}\) is recorded
rather than pooled separately.  Let \(\widetilde H_\Phi^\star\) be the true
projected response and \(\widehat{\widetilde H}_\Phi\) the constructed one, with
\(\|\widehat{\widetilde H}_\Phi-\widetilde H_\Phi^\star\|_2\le\delta_H\).  Then
\(\widetilde Y=\widehat{\widetilde H}_\Phi\mathbf c
+[\widetilde E+(\widetilde H_\Phi^\star-\widehat{\widetilde H}_\Phi)\mathbf c]\),
so response error acts like additional observation error.

\begin{proposition}[Perturbation bound]
\label{prop:operator_perturbation}
Assume \(\widehat{\widetilde H}_\Phi\) has full column rank and
\(\|\widehat{\widetilde H}_\Phi-\widetilde H_\Phi^\star\|_2\le\delta_H\).
The unconstrained least-squares estimate using
\(\widehat{\widetilde H}_\Phi\) satisfies
\begin{equation}
\|\widehat{\mathbf c}-\mathbf c\|_2
\le
\frac{\|\widetilde E\|_2+\delta_H\|\mathbf c\|_2}
{\sigma_J(\widehat{\widetilde H}_\Phi)}.
\label{eq:operator_error_bound}
\end{equation}
\end{proposition}

A short proof, reusing the pseudoinverse bound of
Proposition~\ref{prop:noise_robust}, appears in
Appendix~\ref{app:identifiability_proofs}.
The same norm bound applies to all three discrepancy terms, but their
interpretations are not interchangeable.  Wind and dispersion ensembles
quantify transport uncertainty and may support uncertainty intervals under the
ensemble model.  Alternative inventory maps, locations, labels, or scales are
inventory robustness scenarios, not confidence intervals.  They are reported
separately, and source-specific attribution remains conditional on the supplied
inventory.  By Proposition~\ref{prop:operator_perturbation}, either error type
increases the coefficient error more when \(\sigma_J\) is small.


\subsection{Per-Sensor Source Footprints}
\label{app:footprint_details}

Policy questions are often also spatial, and because the response is linear in
\(\mathbf c\), the fitted solution already localizes each monitor's signal.  For an
observed row \((s,t)\), the background-corrected fit decomposes exactly as
\(\widetilde y_{s,t}=\sum_{k,b}\widetilde H_{\Phi,(s,t),(k,b)}\,\widehat c_{kb}\),
so the contribution of source \(k\) to sensor \(s\) over the record is
\(\widehat Y^{(s)}_k=\sum_t\sum_b
H_{\Phi,(s,t),(k,b)}^{\mathrm{lag}}\widehat c_{kb}\).  The projected form isolates
the identifiable part; the unprojected form gives the raw fitted sensor
contribution before background removal.  The spatial origin of a sensor's signal
is the pullback of its response row onto the source grid: for sensor \(s\) at time
\(t\), \(F_{s,t}(i)=\sum_{\ell\in\mathcal L_t}
[O G_{t,t-\ell}^{\partial}(\widehat{\mathbf w}_{t-\ell:t})]_{s,i}\ge 0\) and
\(F_s(i)=\sum_t F_{s,t}(i)\), a nonnegative field over source cells \(i\)
recording which upwind cells influence sensor \(s\).  Weighting by
\(\phi_b(t-\ell)\), the inventory map \(\mathbf s_k\), and \(\widehat c_{kb}\)
resolves the fitted contribution of each source group to each sensor by cell of
origin.  This is the puff response of Appendix~\ref{subsec:plume_response} read
backward from the sensor and requires no new transport operator.  Per-sensor
attribution uses only the rows of \(\widetilde H_\Phi\) for that sensor: for the
row submatrix \(\widetilde H_\Phi^{(s)}\),
\(\widetilde H_\Phi^{(s)\top}\widetilde H_\Phi^{(s)}\preceq
\widetilde H_\Phi^\top\widetilde H_\Phi\), so deleting rows cannot increase the
smallest singular value or the rank and a single sensor is never more
identifiable than the pooled network.  Footprints and per-sensor contributions
are therefore reported as spatial overlays of the global fit; per-sensor source
\emph{shares} are asserted only at the globally identifiable resolution,
aggregating to a report group wherever the pooled diagnostics require one, which
keeps the spatial view from implying separations the sensing system does not
support.


\subsection{New Delhi Platform Details}
\label{app:platform_details}

This appendix expands the platform construction of
Appendix~\ref{sec:evaluation}: the source inventories
(Table~\ref{tab:new_delhi_inventories}), the kriged initial-condition baseline,
and the auxiliary simulator.

\begin{table}[t]
\centering
\small
\setlength{\tabcolsep}{4pt}
\caption{New Delhi source groups, spatial proxies, and temporal-basis components.}
\label{tab:new_delhi_inventories}
\begin{tabularx}{\columnwidth}{l
  >{\raggedright\arraybackslash}X
  >{\raggedright\arraybackslash}X}
\toprule
Source group & Spatial proxy & Temporal-basis components \\
\midrule
Brick kilns & regional kiln map & alternating 12-hour blocks \\
Industries & regional industry map & continuous + daytime (07--19) \\
Population density & gridded population & peaks 07:00, 13:00, 19:00 \\
Traffic & road-network map & diurnal slots 00/06/12/18 \\
\bottomrule
\end{tabularx}
\end{table}

\noindent\textbf{Observations, sensors, and inventory preparation.}
The hourly government records \citep{cpcb,cpcb_portal} (columns
\texttt{monitor\_id}, \texttt{timestamp\_round}, \texttt{AT}, \texttt{RH},
\texttt{WD}, \texttt{WS}, \texttt{pm10}, \texttt{pm25}) span 21,960 hourly
timestamps from 1 May 2018 through 31 October 2020 in Indian Standard Time;
over all 32 locations, including East Arjun Nagar, which has no readings,
approximately 74.9\% of wind-direction, 72.8\% of wind-speed, and 90.5\% of
PM\(_{2.5}\) entries are observed, with a valid-value mask retained alongside
station coordinates and timestamps.  The two Pusa monitors (\texttt{Pusa\_IMD} and
\texttt{Pusa\_DPCC}) are averaged by timestamp into \texttt{Pusa\_averaged} (mean
coordinates, per-variable hourly averaging), giving the final 32-sensor layout;
PM\(_{2.5}\) is never imputed and its flattened valid-value mask defines
\(M_{\mathcal O}\), with the same ordered sensor--time rows retained in \(Y\),
\(H_\Phi^{\mathrm{lag}}\), \(Q\), and metadata.  Brick-kiln and industry proxies
follow the regional emissions inventory \citep{GUTTIKUNDA2013101}, population
density the gridded population product \citep{ColumbiaUniversity2018}, and the
traffic map is derived from road map data \citep{google_maps}.  Every
\(80\times80\) map is cropped to the New Delhi study window (rows \(21{:}61\),
columns \(16{:}56\)) covering the regulatory network, giving a \(40\times40\) map,
then divided by its own cropped 99th percentile.  Traffic is modeled as a single
road-network source whose nearest-slot maps at 00:00, 06:00, 12:00, and 18:00
define traffic-specific temporal-basis components under the shared dictionary
\(\Phi\) and the fixed-zero mask \(\mathcal F_0\) (assuming a time-invariant
congestion spatial pattern whose amplitude varies); brick-kiln activity uses
deterministic alternating 12-hour blocks, industry a continuous operating fraction
with higher daytime (07:00--19:00) activity, and population a baseline with smooth
cooking peaks.  These are deterministic study assumptions, not observed emissions
truth, and each map's independent proxy normalization is why fitted coefficients
are in normalized proxy units; the group signals are in \(\mu\mathrm g/\mathrm m^3\),
and the shares are not emission shares.

\noindent\textbf{Forward-model baselines and study modes.}
Both forward paths use the New Delhi-derived \(40\times40\) domain and regulatory
sensor coordinates, with the open-boundary Gaussian puff operator evaluated
directly at sensors.  The same operator is the matched generator for controlled
recovery experiments and the inference operator for observed PM\(_{2.5}\); only the
initial-condition baseline differs by mode: a zero-source, zero-initial-state
baseline for controlled operator-matched runs, and a subtracted zero-source
transported kriged-initial-condition baseline for observed runs.  Normal
backgrounds follow Appendix~\ref{subsec:background_construction}, with the base
controlled background the rank-four member.  Where a study is explicitly about
inventory geometry (the coherence, inventory-robustness and missing-source experiments) the real
New Delhi inventory maps are used; the remaining controlled experiments use compact
synthetic source cells for tractability, since their conclusions concern
conditioning, coherence, and recovery rather than specific inventory footprints.

\noindent\textbf{Kriged initial-condition baseline.}
For observed New Delhi runs, the initial pollution field is estimated from the
first hour of the window in which at least three stations report, by a
normalized Gaussian-kernel spatial interpolation on the response grid (a kriging
surrogate).  Its
zero-source transported trajectory, obtained by propagating this initial field as
a single release through the same open-boundary puff operator (its Gaussian kernel
is the advection--diffusion Green's function, so units are preserved), is
subtracted from the sensor observations before fitting source activities.  Source
coefficients are therefore fitted to the residual sensor-time signal after
removing the contribution of the initialized background field and the low-rank
temporal background basis.  Kernel interpolation from sparse first-hour
observations is itself an ill-posed spatial estimate, so this baseline injects
some uncertainty; the effect is limited by declaring and subtracting the baseline
before fitting and by the low-rank background basis \(Q\), which absorbs residual
broad structure.  Controlled operator-matched experiments may instead use a zero
initial field when the synthetic data are generated from the same
zero-initial-state response.

\noindent\textbf{Auxiliary advection--diffusion simulator.}
The auxiliary advection--diffusion simulator uses first-order upwind advection, a
five-point Laplacian, a two-stage Heun update, diffusivity \(3\times10^{-4}\), and
edge-hold boundaries.  It provides a structural forward-model mismatch: data
generated by this distinct transport operator are fit with the open-boundary puff
response, testing the puff approximation against a different operator family.
This mismatch also drives the residual-adequacy test of
Section~\ref{sec:identifiability_theory}, letting us check its power against structural
transport misspecification and not only against missing sources.


\subsection{Wind-Imputation Validation}
\label{subsec:results_wind_imputation}
The learned wind model did not beat a non-spatial city-mean baseline on held-out
station--time vectors, so the spatially resolved kernel imputer is used throughout (the city mean lacks the
spatial resolution the response operator requires).  Dense wind truth is
unavailable, so gridded-field accuracy is assessed only through the controlled
synthetic-wind experiments, where the true wind is known.

